% Plan de tesis de licienciatura Diciembre 2005.
% Pablo Listingart

\documentclass[a4paper, 12pt]{article}
\usepackage[spanish]{babel}
%\usepackage{makeidx}
\usepackage[pdftex]{graphicx}

\addtolength{\topmargin}{-2cm}
\addtolength{\textheight}{4cm}
\makeindex 

\begin{document}
\DeclareGraphicsExtensions{.jpg,.pdf,.mps,.png}

\title {
\begin{figure}[h]
\begin{center}
\includegraphics[keepaspectratio, width=3cm]{escudo}
\end{center}
\end{figure}
Universidad de Buenos Aires\\
Facultad de Ciencias Exactas y Naturales\\
Departamento de Computaci\'on\\
\vspace{10mm}
\Huge\textbf{Computaci\'on de alto rendimiento en el estudio de reacciones qu\'imicas en soluciones 
a trav\'es de m\'etodos cl\'asicos y cu\'anticos.}
\vspace{10mm}
}

\author{por \\
	Pablo Listingart\\
        \\
        Director de Tesis\\
        Dr. Guillermo Marshall\\
        \\
        Propuesta de tesis para optar al grado de \\
        Licenciado en Ciencias de la Computaci\'on\\
        \\
        }
\date{Diciembre de 2005}

\maketitle

\pagebreak

\section*{Propuesta}

Construcci\'on de un software paralelo de simulaci\'on num\'erica a partir de una versi\'on serial preexistente y su implementaci\'on en un cluster Beowulf para el estudio y predicci\'on de reacciones qu\'imicas en soluci\'on. 

La combinaci\'on del hardware y software mencionado y el an\'alisis de distintas estrategias de paralelizaci\'on (incluyendo diversas modalidades de pasaje de mensajes) permitir\'an la resoluci\'on de problemas m\'as realistas y previamente inabordables, obteniendo un algoritmo escalable.


\section*{Objetivos}

\begin{itemize}
\item Desarrollar un programa paralelo que resuelva num\'ericamente un esquema h\'ibrido cu\'antico-cl\'asico (QM-MM) en el cual el solvente es tratado cl\'asicamente.
\item Comparar los resultados num\'ericos obtenidos frente a los obtenidos en forma serial del mismo problema.
\item Incorporar herramientas de visualizaci\'on para realizar el an\'alisis de 
los datos obtenidos.
\item Aplicar el programa obtenido a tama\~nos de problemas m\'as realistas y previamente inabordables.
\end{itemize}

\section*{Elementos a utilizar}

\begin{itemize}
\item Sistema Operativo: Linux, Windows.
\item Lenguaje de programaci\'on: Fortran 77.
\item M\'aquina con suficiente memoria y poder de c\'alculo para poder 
realizar simulaciones que puedan aproximarse a la realidad f\'isica.
\item C\'odigo fuente de la versi�n serial que soluciona el modelo.
\end{itemize}

\section*{Plan inicial}

\begin{enumerate}
\item Interiorizarse de la realidad qu\'imica que ocurre durante un 
experimento en soluciones a trav\'es de m\'etodos cl\'asicos y cu\'anticos.
\item Analizar el programa serial utilizado para la resoluci\'on de sistemas qu\'imicos.
\item Determinar las t\'ecnicas posibles de paralelizaci\'on
\item Realizar la pralelizaci\'on del programa serial.
\item Realizar experimentos para el estudio de reacciones qu\'imicas en soluciones a trav\'es de m\'etodos cl\'asicos y cu\'anticos.
\item Analizar los resultados y relizar pruebas de performance.
 \end{enumerate}

\section*{Bibliograf\'ia}

   % \begin{thebibliography}{9}
   % 
   % \bibitem 
   % 
   % \bibitem 
   % 
   % \bibitem
   % 
   % \bibitem
   % 
   % \bibitem
   % 
   % \bibitem
   % 
   % \end{thebibliography}

\end{document}
